\section*{Capítulo \ref{ch:b2:vertebrate-development} --- Desarrollo embrionario de los vertebrados}

\begin{solution}{exo:b2:vertebrate-development:1}
\omterm{def:b2:vertebrate-development:egg}{Segmentación}: mitosis rápidas sin crecimiento, que producen la \omterm{def:b2:vertebrate-development:blastula}{blástula}.
\omterm{def:b2:vertebrate-development:blastula}{Blástula}: una esfera hueca o un disco de células pequeñas alrededor del
\omterm{def:b2:vertebrate-development:blastula}{blastocele}. Gastrulación: los movimientos celulares que llevan el
mesodermo y el endodermo al interior y producen la gástrula de tres capas,
con una cavidad intestinal y una notocorda. Neurulación: el plegamiento
del ectodermo dorsal en el \omterm{prop:b2:vertebrate-development:neurulation}{tubo neural}, con la \omterm{def:b2:vertebrate-development:egg}{segmentación} de los
somitos, que produce la néurula con el plan corporal de los vertebrados.
\end{solution}

\begin{solution}{exo:b2:vertebrate-development:2}
Ectodermo: epidermis, encéfalo y médula espinal, nervios periféricos
(\omterm{prop:b2:vertebrate-development:neurulation}{cresta neural}), cristalino y células pigmentarias. Mesodermo: notocorda,
músculo, hueso y cartílago, riñón, corazón y sangre, dermis. Endodermo:
el revestimiento del intestino, hígado, páncreas, pulmones, tiroides.
\end{solution}

\begin{solution}{exo:b2:vertebrate-development:3}
Poco \omterm{def:b2:vertebrate-development:egg}{vitelo}: \omterm{def:b2:vertebrate-development:egg}{segmentación} completa y casi igual, y gastrulación por
involución a través de un blastoporo redondo en una esfera hueca (rana:
completa pero desigual, con las células vegetativas cargadas de \omterm{def:b2:vertebrate-development:egg}{vitelo}
grandes y lentas). Mucho \omterm{def:b2:vertebrate-development:egg}{vitelo} (pollo): \omterm{def:b2:vertebrate-development:egg}{segmentación} limitada a un disco
situado sobre el \omterm{def:b2:vertebrate-development:egg}{vitelo}, y gastrulación a través de una hendidura, la
línea primitiva, en el disco plano, con las capas extendiéndose sobre el
\omterm{def:b2:vertebrate-development:egg}{vitelo} en lugar de encerrar una cavidad.
\end{solution}

\begin{solution}{exo:b2:vertebrate-development:4}
La \omterm{def:b2:vertebrate-development:induction}{inducción} es el cambio de destino de un tejido competente por una señal
procedente de un tejido vecino. Las células vegetativas inducen mesodermo
en la zona marginal situada encima; el \omterm{def:b2:vertebrate-development:induction}{organizador} (mesodermo del labio
dorsal) induce la placa neural en el ectodermo dorsal; la vesícula óptica
induce el cristalino en el ectodermo de la cabeza.
\end{solution}

\begin{solution}{exo:b2:vertebrate-development:5}
$2^{12} = 4096$ células de radio $1/2^{4} = \qty{62.5}{\micro m}$;
superficie multiplicada por $2^{4} = 16$. La superficie adicional es la
membrana de los epitelios que la gastrulación plegará en piel e
intestino, y el área a través de la cual la \omterm{def:b2:vertebrate-development:blastula}{blástula} intercambia con su
entorno.
\end{solution}

\begin{solution}{exo:b2:vertebrate-development:6}
$10^{-4}\times 2^{k} = 0.4$: $2^{k} = 4000$, $k = 12$. Con cuatro veces
el ADN, la relación parte de $4\times 10^{-4}$ y alcanza $0.4$
con $2^{k} = 1000$: la décima división, dos antes.
\end{solution}

\begin{solution}{exo:b2:vertebrate-development:7}
$k = \ln 2/600 = \qty{1.16e-3}{s^{-1}}$; $\lambda = \sqrt{2\times
10^{-12}/1.16\times 10^{-3}} = \qty{42}{\micro m}$. Umbrales en
$x = \lambda\ln 2 = \qty{29}{\micro m}$ y $\lambda\ln 10 =
\qty{96}{\micro m}$: bandas de 29, 67 y el resto.
\end{solution}

\begin{solution}{exo:b2:vertebrate-development:8}
Duplicar la fuente desplaza ambas fronteras hacia fuera en $\lambda\ln
2 = \qty{29}{\micro m}$: la primera banda duplica su anchura y la segunda
la conserva. Duplicar $k$ divide $\lambda$ por $\sqrt 2$
(\qty{29}{\micro m}) y, como $c_0 = j/\sqrt{Dk}$, divide también $c_0$ por
$\sqrt 2$: las fronteras pasan a $\lambda'\ln(c_0'/T) =
29\times(0.69 - 0.35) = \qty{10}{\micro m}$ y $29\times(2.30 -
0.35) = \qty{57}{\micro m}$ --- ambas bandas se estrechan.
\end{solution}

\begin{solution}{exo:b2:vertebrate-development:9}
Se injertó un labio dorsal de una gástrula de tritón no pigmentada en el
lado ventral de otra pigmentada; se enrolló hacia dentro y el huésped
formó un segundo embrión --- \omterm{prop:b2:vertebrate-development:neurulation}{tubo neural}, somitos, intestino --- unido al
primero. La diferencia de pigmento mostró que el segundo eje estaba
formado por células del huésped, de modo que el injerto las había
inducido en lugar de construir el eje por sí mismo. Descartó la
autodiferenciación del injerto como explicación, y mostró que una región
pequeña lleva la señal que organiza todo el eje.
\end{solution}

\begin{solution}{exo:b2:vertebrate-development:10}
(a) Piel ventral: el ectodermo temprano es competente pero aún no ha sido
inducido, de modo que sigue a su nuevo entorno. (b) Tejido nervioso: en la
gástrula tardía ya ha sido inducido y está determinado. (c) Placa neural:
el ectodermo ventral de la gástrula tardía sigue siendo competente, y la
notocorda lo induce. (d) Ninguna estructura dorsal: una bola de tejido
ventral, puesto que nada induce los destinos dorsales en el mesodermo ni
los neurales en el ectodermo.
\end{solution}

\begin{solution}{exo:b2:vertebrate-development:11}
A \qty{18}{\celsius}: $90 + 11\times 30 = \qty{420}{min}$, 7 horas.
A \qty{10}{\celsius}: \qty{17.5}{h}. La transición cae en la duodécima
división en ambos casos porque lo que la desencadena es la relación entre
ADN y citoplasma, que depende del número de divisiones y no de lo que
hayan tardado: el reloj cuenta duplicaciones, no horas.
\end{solution}

\begin{solution}{exo:b2:vertebrate-development:12}
En el desarrollo regulativo las células adquieren su destino por
\omterm{def:b2:vertebrate-development:induction}{inducción} de sus vecinas, de modo que el embrión puede reconstruirse tras
una pérdida (gemelos a partir de un huevo dividido, regeneración de una
región injertada); en el desarrollo en mosaico los destinos se heredan con
determinantes citoplasmáticos localizados, de modo que cada blastómero
solo forma su parte y un embrión dividido da dos medias larvas. La
diferencia es de grado: la media luna gris de la rana es un determinante,
y los embriones en mosaico recurren más tarde a inducciones; todo embrión
mezcla herencia y conversación, y los vertebrados se apoyan en la
conversación durante más tiempo.
\end{solution}

\begin{solution}{pb:b2:vertebrate-development:1}
\textbf{1.} Huevo $\tfrac43\pi(0.7)^{3} = \qty{1.44}{mm^3}$; núcleo
$\tfrac43\pi(0.03)^{3} = \qty{1.1e-4}{mm^3}$; relación $8\times
10^{-5}$.
\textbf{2.} $2^{k} = 4000$: $k = 12$, 4096 células.
\textbf{3.} $75 + 11\times 30 = \qty{405}{min}$, unas 6,75 horas.
\textbf{4.} Radio $0.7/2^{4} = \qty{44}{\micro m}$; relación $8\times
10^{-5}\times 4096 = 0.32$ --- la de una célula ordinaria: los núcleos han
alcanzado al citoplasma.
\textbf{5.} $2^{4} = 16$.
\textbf{6.} Unas $4095\,c$ de ADN; el huevo contiene suficientes histonas,
polimerasas y nucleótidos almacenados para doce divisiones, de modo que no
hace falta transcripción.
\textbf{7.} Se activan los genes cigóticos, las divisiones se ralentizan y
pierden la sincronía, y las células se vuelven móviles. Inyectar ADN
adicional adelanta la transición tantas divisiones como veces está el ADN
en exceso, algo que un contador de divisiones o de tiempo no podría
explicar.
\textbf{8.} Frente al punto de entrada del espermatozoide, en la media
luna gris fijada por la rotación de la corteza del huevo tras la
fecundación.
\textbf{9.} Primero el mesodermo dorsal (placa precordal y después
notocorda), hacia el techo del arquénteron en la línea media; el
mesodermo lateral (somitos, lámina lateral) a su lado; el endodermo
reviste el arquénteron; el ectodermo se extiende por el exterior por
epibolia.
\textbf{10.} $\qty{2000}{\micro m}/\qty{600}{min} = \qty{3.3}{\micro
m/min}$: tres veces un fibroblasto --- una lámina coordinada avanza más
deprisa que una célula aislada.
\textbf{11.} Los movimientos celulares necesitan productos de los genes
cigóticos (nuevas moléculas de adhesión, motilidad) y ciclos ralentizados;
antes de la transición las células solo pueden dividirse.
\textbf{12.} Todo cordado tiene notocorda en algún estadio, y los
vertebrados construyen la columna vertebral a su alrededor; el \omterm{prop:b2:vertebrate-development:neurulation}{tubo neural}
es inducido por ella y es la consecuencia, no el origen, del plan.
\textbf{13.} Sin la contracción de las células en botella no se forma el
blastoporo, nada entra, y el embrión sigue siendo una bola con sus capas
en su sitio: sin intestino, sin notocorda, sin eje.
\textbf{14.} $k = \ln 2/1200 = \qty{5.8e-4}{s^{-1}}$; $\lambda =
\sqrt{10^{-12}/5.8\times 10^{-4}} = \qty{42}{\micro m}$.
\textbf{15.} $x_1 = 42\ln 2.5 = \qty{38}{\micro m}$; $x_2 = 42\ln 12.5
= \qty{105}{\micro m}$.
\textbf{16.} Bandas de \qty{38}{\micro m} (2,5 células), \qty{67}{\micro
m} (4,5 células) y el resto.
\textbf{17.} $x_1 = 42\ln 1.25 = \qty{9}{\micro m}$; $x_2 = 42\ln 6.25
= \qty{77}{\micro m}$: ambas fronteras se desplazan hacia dentro en
$\lambda\ln 2 =
\qty{29}{\micro m}$; la primera banda se reduce de 38 a 9, la banda media
conserva su anchura y la externa crece.
\textbf{18.} Trasladada antes de la lectura, adopta el destino A --- lee
la concentración del lugar donde está ahora; trasladada después, conserva
el destino B --- ya estaba determinada.
\textbf{19.} $x_T = \lambda\ln(c_0/T)$: un cambio de la fuente en un factor
$f$ desplaza cada frontera en $\lambda\ln f$, de modo que un error del
doble solo desplaza el patrón en $0.7\lambda$, una fracción de un campo que
abarca varios $\lambda$ --- el patrón es casi insensible a la cantidad
exacta de morfógeno.
\textbf{20.} Labio dorsal en posición ventral: un segundo eje. Zona
marginal ventral en posición dorsal: su entorno la vuelve a especificar y
forma tejido dorsal; nada adicional.
\textbf{21.} En una néurula el ectodermo ya no es competente: el injerto
forma notocorda, pero no induce ningún segundo \omterm{prop:b2:vertebrate-development:neurulation}{tubo neural}.
\textbf{22.} Dividido en el plano de la media luna: dos embriones
completos; perpendicularmente: un embrión y una bola de tejido ventral.
\textbf{23.} Solo: una bola epidérmica. Con células vegetativas: mesodermo
--- músculo, notocorda --- inducido en el casquete.
\textbf{24.} Para demostrar que el segundo eje se había inducido en
células del huésped y no lo había construido el injerto; si el injerto
hubiera sido indistinguible, no podría haberse descartado la
autodiferenciación. Un injerto de otra especie funciona (usaron dos
especies de tritón), y eso es lo que hizo distinguibles las células.
\textbf{25.} Duodécima \omterm{def:b2:vertebrate-development:egg}{segmentación}, 4096 células, a las 6,75 horas
aproximadamente; bandas de 38 y \qty{67}{\micro m}, y todo lo que queda
más allá de \qty{105}{\micro
m}.
\end{solution}
